\chapter{Riproducibilità della valutazione}
\label{app:riproducibilita}

L'intera valutazione presentata in \Cref{ch:valutazione} è riproducibile a partire dal repository pubblico, in due modalità di livello crescente di rigore. La modalità leggera consiste nell'eseguire il solo rule engine sull'insieme dei casi di riferimento, che richiede pochi secondi su un portatile e non dipende dall'LLM locale; la modalità completa consiste nell'eseguire l'intera pipeline overnight, che richiede una GPU NVIDIA con almeno 8\,GB di memoria e tempi dell'ordine delle quattordici-sedici ore. Le versioni dei componenti software fissati per la \cleanroomrun\ sono Python 3.12 per il rule engine e per le metriche di valutazione, contenitore CUDA Ubuntu 22.04 con Python 3.12 per la pipeline RAG, Ollama come runtime dell'LLM e il modello \llmmodel\ scaricato in formato GGUF Q4\_K\_M. Il modello di embedding è \embedmodel, il reranker è \reranker\ e per le metriche di fedeltà NLI sono stati impiegati due checkpoint mDeBERTa: \texttt{mDeBERTa-v3-base-mnli-xnli} (metrica \emph{nli\_faithfulness}) e la sua evoluzione \nlimodel\ (metrica \emph{semantic\_faithfulness\_v2}); il confronto fra i due, a parità di metrica, è discusso in \Cref{sec:rq3}.

La modalità leggera si attiva con il comando di seguito riportato, da eseguire nella radice del repository dopo aver attivato l'ambiente virtuale del rule engine.

\begin{verbatim}
cd Note_AIFA/aifa_rule_engine
source .venv/bin/activate
python -m pytest tests/ -q                     # 435 test, 0,7 s
cd ..                                          # radice Note_AIFA
python -m evaluation.scripts.evaluate_rule_engine \
       --json-report evaluation/results/rule_engine_report.json
\end{verbatim}

L'output atteso è una Macro F1 pari a $1{,}0000$ su \goldn\ casi e una distribuzione delle classi conforme a quanto riportato in \Cref{tab:rule-engine}. La pipeline completa, comprensiva di indicizzazione delle Note, valutazione del retrieval, generazione delle spiegazioni LLM, calcolo delle metriche di fedeltà e produzione dei \emph{per-case report}, si attiva mediante lo script orchestratore di seguito riportato.

{\footnotesize
\begin{verbatim}
cd Note_AIFA
bash tools/run_full_overnight.sh                 # ~14-16 ore
bash tools/run_full_overnight.sh --resume        # ripresa post-interruzione
bash tools/run_full_overnight.sh --skip-ragas    # esclude RAGAS (~3 ore)
\end{verbatim}
}

Lo script suddivide la pipeline in ventitré stage discreti, ciascuno con un checkpoint persistente in \texttt{evaluation/results/.checkpoints/}; il flag \texttt{--resume} riprende l'esecuzione dal primo stage non completato e consente di superare riavvii o sospensioni della macchina senza dover ripartire da zero. Lo stage di RAGAS, il più lungo a causa dell'inferenza ripetuta dell'LLM giudice, è opzionalmente disattivabile con \texttt{--skip-ragas}; in questo caso le metriche faithfulness e answer relevancy sono assenti dal report finale, mentre tutte le altre rimangono prodotte.

Per garantire l'indipendenza dei risultati dallo stato precedente del sistema, esiste una modalità \emph{cleanroom} di livello superiore, che cancella esplicitamente lo store ChromaDB, la cartella dei risultati e tutti i checkpoint prima di rieseguire la pipeline. Si attiva con lo script di seguito riportato e produce, alla conclusione, uno snapshot con timestamp in \texttt{audit/.snapshots/}, utile per il confronto tra run successive.

{\footnotesize
\begin{verbatim}
cd Note_AIFA
bash tools/full_cleanroom_v2.sh                  # wipe + ingest + tests + eval
\end{verbatim}
}

\begin{sloppypar}
Tutti i risultati prodotti dalla \cleanroomrun\ sono salvati in \texttt{Note\_AIFA/\allowbreak evaluation/\allowbreak results/} sotto forma di file JSON canonici, accompagnati da un riepilogo testuale in \texttt{OVERNIGHT\_\allowbreak SUMMARY.md}. La verifica della consistenza fra i numeri citati nella tesi e quelli effettivamente presenti nei JSON è automatizzata da un piccolo script di check in \texttt{tools/\allowbreak verify\_\allowbreak thesis\_\allowbreak numbers.py}, che il revisore clinico può eseguire per confermare l'allineamento.
\end{sloppypar}

Sul piano della validazione clinica dei casi di riferimento, un sottoinsieme di casi-frontiera è stato sottoposto a revisione esterna da parte di un medico di medicina generale con esperienza prescrittiva sulle quattro Note codificate. La revisione ha riguardato venti casi selezionati per soglie cliniche borderline (LDL al confine di $100$~mg/dL per \notetredici, CHA$_2$DS$_2$-VASc al confine $2{-}3$ per \notenova, doppia indicazione gastroprotettiva per \notezero, doppio fattore di rischio gastrointestinale per \notesessanta). Il pacchetto di revisione è documentato nel file \texttt{audit/CLINICAL\_REVIEW\_PACKET\_2026-05-13.pdf} ed è disponibile insieme al repository.
